\section{Introduction}
\label{sec:introduction}
\subsection{Categorized literature review}
\label{sec:literature}

\textbf{Prediction-powered inference.} Prediction-powered inference uses a large set of machine predictions together with a smaller labeled sample, correcting the prediction-based estimate with labeled residuals so that the statistical target does not require the predictor to be correct \citep{angelopoulos2023ppi}. This separation between predictive quality and inferential validity is the starting point for the present construction. Randomized observations play the role of causally valid labeled data, while observational outcome models supply potentially useful predictions.

\textbf{Adaptive prediction-powered inference.} PPI++ introduces data-adaptive power tuning that adjusts how strongly predictions enter an estimating equation and reports computational and statistical gains over the original construction \citep{angelopoulos2024ppipp}. Its central lesson is that a prediction should be used in proportion to its empirical value rather than being accepted as truth. Our coefficient $\omega$ follows this control-variate logic, but the target and score are causal quantities identified by the randomized design.

\textbf{Causal data fusion and transport.} General causal data-fusion theory distinguishes data collected under different regimes and populations and asks whether a target causal query is identified from their combination \citep{bareinboim2016fusion}. Reviews of causal data integration likewise separate randomized trials, observational studies, historical controls, and heterogeneous target populations \citep{shi2022integration}. The present paper assumes that the trial already identifies the target effect and studies an estimation question: how can observational predictions improve precision without replacing randomized identifying information?

\textbf{Combining randomized and observational estimates.} Existing estimators combine experimental and observational effect estimates through shrinkage, bias detection, or adaptive weighting. Shrinkage estimators can reduce risk when observational bias is sufficiently controlled \citep{rosenman2020shrinkage}, while adaptive combination methods move toward the randomized estimator when observational bias is detected \citep{cheng2021adaptive}. These approaches combine effect estimates from the two sources. Our construction instead inserts the observational outcome model into an RCT-valid pseudo-outcome and separately centers a prediction-based covariate statistic.

\textbf{CATE learners and orthogonal losses.} Meta-learners convert standard regression algorithms into conditional effect estimators \citep{kunzel2019metalearners}. The R-learner residualizes outcomes and treatments to isolate treatment-effect heterogeneity \citep{nie2020rlearner}, while doubly robust pseudo-outcome regression can retain favorable error structure under nuisance estimation \citep{kennedy2022drlearner}. Direct multi-source learners have also been proposed for heterogeneous causal data fusion \citep{li2022directfusion}. These results motivate treating the prediction-powered construction as a loss wrapper that can be paired with a linear sieve or a fixed learned representation.

\textbf{Covariate shift and balancing.} When trial and observational covariate marginals differ, exact density-ratio weighting transports fixed covariate functions between populations. Riesz-representer methods provide an optimization perspective on such linear functionals and on debiasing \citep{chernozhukov2021riesz}. Density-ratio estimation, generic balancing, and Riesz learning are established tools. The project-specific question is whether they can be targeted to the generated CATE loss class $\{(\widehat g-t)^2:t\in\mathcal T\}$, where $\widehat g$ is the fitted OBS treatment-effect prediction defined later in Equation~\eqref{eq:obs-effect-prediction}, while preserving an RCT-defined risk; we treat this specialization as a working research direction rather than a settled novelty claim.
